\section{Work Completed and Technical Details}

\subsection{Done so far}

We downloaded the published release once and kept it as a local copy, so later updates to the public collection stay out of the analysis, and we analyze its tables with Python and Polars.
Names are extracted with GLiNER,\footnote{\url{https://arxiv.org/abs/2311.08526}} and the finished work covers the items below.

\begin{itemize}
  \item A frozen local copy of the collection.
  \item Checks of star counts, host-repository license labels, and the text stored beside a skill.
  \item A label-only model run on 500 companion texts.
\end{itemize}

Every notebook sets one shared random seed when it loads the tables, and the 500-text draw uses that seed, so a rerun returns the same texts, each still tied to the file it came from.
The twelve-text trial in the last section drew its sample with its own fixed seed.
Dated PDF exports of the explore notebooks are kept in \path{scripts/exports}, so their outputs can be read without running the notebooks, which may still change.

\subsection{Score cutoff and accuracy}

The model gives each hit a score, and a cutoff on that score decides which hits we keep.
We are still studying how that cutoff relates to accuracy, and for now we test the simplest pairing, so a cutoff of~$0.8$ is checked against an accuracy near~$80\%$, although that pairing is only an arbitrary starting point.
To measure accuracy, the authors will label the taught tools in a random sample of the population, drawn with the shared seed and sized before anyone sees the model output.
We will report precision, the share of extracted tools that are truly taught tools, and recall, the share of truly taught tools the model finds, along with inter-rater agreement on a subset that two authors label independently, but no such number appears here because that sample has not been labeled.
